% Generated draft: requires mathematical and source review.
\documentclass[11pt]{article}
\usepackage[margin=25mm]{geometry}
\usepackage{fontspec}
\setmainfont{Latin Modern Roman}
\usepackage{amsmath,amssymb,mathtools,unicode-math}
\setmathfont{Latin Modern Math}
\usepackage{enumitem,needspace,xurl,hyperref}
\setlength{\parindent}{0pt}
\setlength{\parskip}{5pt}
\setlength{\emergencystretch}{3em}
\newcommand{\N}{\mathbb{N}}
\newcommand{\Z}{\mathbb{Z}}
\newcommand{\Q}{\mathbb{Q}}
\newcommand{\R}{\mathbb{R}}
\newcommand{\C}{\mathbb{C}}
\newcommand{\F}{\mathbb{F}}
\newcommand{\A}{\mathbb{A}}
\newcommand{\PP}{\mathbb{P}}
\newcommand{\E}{\mathbb{E}}
\newcommand{\eps}{\varepsilon}
\newcommand{\dd}{\,\mathrm{d}}
\begin{document}
\section*{UnsolvedMath ID: 20002305}
\textbf{Source problem code:} catalogue entry
\subsection{Definitions and mathematical statement}
Question 3 asks how difficult it is to find a presentation of \(G\) in terms of \(M\). The excerpt omits the preceding definitions of \(G,M\). The supplied verification uses the following questions, including a three-to-one map from a strongly minimal group to \((\Q,+)\), to identify computable-model-theoretic presentations and finite covers/quotients as the context; it does not reproduce a complete hypothesis set for Question 3.

Let \(D(M)\) denote the atomic diagram of a countable presentation of \(M\). For the fixed-interpretation reading, suppose an effective interpretation of a group \(G\) in \(M\), including its parameter tuple and defining scheme, is supplied. The aid reports a presentation \(\widetilde G\cong G\) with
\[
 D(\widetilde G)\le_T D(M),
\]
without a jump. This uses the cited \(\Delta_1\)/effective-interpretation setting and is not an algorithm for discovering the interpretation from an arbitrary presentation of \(M\).

The candidate strengthens the quotient-data hypothesis in one finite-fiber case. Relative to \(A=D(M)\), suppose:
\begin{enumerate}
\item \(U\) is an \(A\)-c.e. domain of finite tuples, with decidable equality of tuple codes;
\item \(E\) is an \(A\)-c.e. equivalence relation on \(U\), and every class has the same supplied finite size \(k\ge1\);
\item an \(A\)-c.e. ternary graph on \(U\) induces a total, single-valued multiplication on \(U/E\), making it an infinite group.
\end{enumerate}
``Single-valued'' is at the quotient level: output representatives may differ within one class. The aid reports that \(E\) is \(A\)-decidable on inputs promised to lie in \(U\), a representative presentation on \(\omega\) can be built with group operations computable in \(A\), and an \(A\)-computable quotient map \(q:U\to\omega\) is available with
\[
 q(x)=q(y)\quad\Longleftrightarrow\quad E(x,y).
\]
The map is computable on its represented domain; no decidability of membership in the c.e. domain \(U\) is asserted. Knowing the common class size is essential to the stated criterion.

The aid contrasts this with c.e. equivalence relations whose class sizes vary between one and two: there are such relations with no computable equality-reflecting quotient map and no computable complete transversal. A complete transversal meets every class exactly once; this does not exclude computable partial sets of pairwise inequivalent elements. The contrast concerns quotient maps/representatives, not an automatic noncomputability theorem for every abstract group using such data.

The unresolved task in the aid is uniformly discovering and certifying an interpretation scheme and its parameters from an arbitrary copy of \(M\). No discovery procedure, sharp complexity bound for that task, or full solution of the contextual original question is supplied.
\subsection{Short English statement}
The source asks for the computable cost of presenting an associated group but omits its earlier definitions. The aid separates using a supplied effective interpretation from discovering one. Its known constant-fiber criterion computes quotient equality, a presentation, and a quotient map without a jump, while variable small fibers can obstruct complete representative selection.
\subsection{Sources}
\begin{itemize}
\item Local UnsolvedMath registry, permanent ID 20002305, exact source code catalogue entry; canonical \texttt{aim-logic-notes.json}, \texttt{notes[80]}, tag \texttt{question}; AIM Workshop Problem Lists, Logic, \emph{Computable stability theory}, Question 3. Created and updated 14 August 2026.
\item Source URL \url{https://aimath.org/pastworkshops/computestabproblems.pdf}. Supplied verification draws contextual evidence from the next two questions and finds no OCR corruption. The PDF was not accessed.
\item Andrews--Medvedev's \(\Delta_1\)-interpretation statement and a later effective-interpretation framework are named by the aid, without full titles, theorem numbers, publication details, or separate URLs.
\item Attempt 1: \texttt{partial\_result}, problem status at run \texttt{partially\_solved}, rigor label \texttt{proved}, novelty confidence \texttt{low}; classification \texttt{PARTIAL-PROGRESS}, registry status \texttt{partially\_solved}. Candidate: constant-fiber no-jump quotient criterion. These are dated source labels, not acceptance.
\item Ulam AI / UnsolvedMath Contributors; CC BY 4.0. Machine-generated research aid, not peer reviewed; independently verify before citation.
\end{itemize}
\end{document}
